% Part: turing-machines 
% Chapter: undecidability
% Section: introduction

\documentclass[../../../include/open-logic-section]{subfiles}

\begin{document}

\olfileid{tur}{und}{int} 
\olsection{ভূমিকা}

প্রতিটি অপেক্ষক, এমনকি প্রতিটি পাটিগাণিতিক অপেক্ষকও গণনসাধ্য হতে পারে
না—এ কথা আপাতদৃষ্টিতে স্পষ্ট মনে হতে পারে। অপেক্ষক এত বেশি, আর তাদের আচরণও
এত জটিল। যেমন, তেজস্ক্রিয় কণার ক্ষয় কিংবা অন্য কোনো বিশৃঙ্খল বা যদৃচ্ছ
আচরণ থেকে সংজ্ঞায়িত অপেক্ষক। ধরা যাক, কোনো নির্দিষ্ট মুহূর্ত থেকে এক
সেকেন্ডের ব্যবধানগুলি গুনতে শুরু করে $f(n)$ অপেক্ষকটিকে ওই আরম্ভিক
মুহূর্তের পর $n$-তম এক-সেকেন্ডের ব্যবধানে মহাবিশ্বে যতগুলি কণা ক্ষয় হয়,
সেই সংখ্যা হিসেবে সংজ্ঞায়িত করলাম। যে অপেক্ষক কখনও গণনা করার আশা করা যায়
না, তার সম্ভাব্য উদাহরণ এটি বলে মনে হয়।

কিন্তু এমন অপেক্ষক কীভাবে গণনা করা যায় তা কল্পনা করতে না-পারা এক কথা, আর
সেগুলি অগণনসাধ্য—এ কথা সত্যিই প্রমাণ করা সম্পূর্ণ অন্য কথা। বাস্তবে যে
অপেক্ষকগুলি আশাতীত জটিল বলে মনে হয়, বিমূর্ত অর্থে তারাও গণনসাধ্য হতে পারে।
যেমন, ধরা যাক মহাবিশ্বটি সময়ের দিক থেকে সসীম—কোনো কোনো মহাজাগতিক তত্ত্বের
পূর্বানুমান অনুযায়ী সুদূর ভবিষ্যতে একদিন মহাবিশ্ব সংকুচিত হয়ে একটিমাত্র
বিন্দুতে পরিণত হবে। তাহলে ওই আরম্ভিক মুহূর্ত থেকে কেবল সসীমসংখ্যক (যদিও
অবিশ্বাস্য রকম বড়) সেকেন্ডের জন্যই $f(n)$ সংজ্ঞায়িত। আর কেবল সসীমসংখ্যক
নিবেশের জন্য সংজ্ঞায়িত যেকোনো অপেক্ষক গণনসাধ্য: আমরা একটি বিরাট সারণিতে
সব নির্গম লিখে দিতে পারি, অথবা একটি অত্যন্ত বড় টুরিং যন্ত্রের
দশা-অবস্থান্তর চিত্রে সেগুলির সংকেতায়ন করতে পারি।

যেসব অপেক্ষকের মান হ্যাঁ/না প্রশ্নের উত্তর দেয়, তাদের বিশেষ ঘটনাগুলি
নিয়ে আমরা প্রায়ই আগ্রহী হই। যেমন, “$n$ কি একটি মৌলিক সংখ্যা?” প্রশ্নটির
সঙ্গে নিচের অপেক্ষকটি যুক্ত:
\[
\fn{isprime}(n) = \begin{cases}
  1 & \text{যদি $n$ মৌলিক হয়}\\
  0 & \text{অন্যথায়।}
  \end{cases}
\]
কোনো হ্যাঁ/না প্রশ্নের সঙ্গে যুক্ত $1/0$-মানবিশিষ্ট অপেক্ষকটি কার্যকরভাবে
গণনসাধ্য হলে আমরা বলি, প্রশ্নটি \emph{কার্যকরভাবে নির্ণয় করা যায়}।

কার্যকরভাবে গণনা করা যায় না এমন অপেক্ষক, অথবা কার্যকরভাবে নির্ণয় করা
যায় না এমন সমস্যা যে আছে, তা গাণিতিকভাবে প্রমাণ করতে গণনার একটি নির্দিষ্ট
মডেল স্থির করা এবং এমন অপেক্ষক বা সমস্যা দেখানো অপরিহার্য, যা ওই মডেল
গণনা বা নির্ণয় করতে পারে না। যেমন, আমরা দেখাতে পারি যে প্রতিটি অপেক্ষক
টুরিং যন্ত্রে গণনা করা যায় না এবং প্রতিটি সমস্যা টুরিং যন্ত্রে নির্ণয় করা
যায় না। তখন চার্চ--টুরিং থিসিস আহ্বান করে সিদ্ধান্ত নিতে পারি যে কেবল
টুরিং যন্ত্রই প্রতিটি অপেক্ষক গণনা করার মতো শক্তিশালী নয় তা নয়, কোনো
কার্যকর পদ্ধতিই তা পারে না।

এমন নেতিবাচক ফল প্রমাণের চাবিকাঠি হলো, টুরিং যন্ত্রগুলিকেই সংখ্যা দেওয়া
যায়। এটি করার সবচেয়ে সহজ উপায় হলো সেগুলিকে তালিকায়িত করা—সম্ভবত টুরিং
যন্ত্র ও তাদের প্রোগ্রাম লেখার একটি নির্দিষ্ট পদ্ধতি স্থির করে, তারপর
সুশৃঙ্খলভাবে তাদের তালিকা বানিয়ে। এটি যে করা যায় তা বুঝলেই সহজ মাত্রা-
বিবেচনা থেকে টুরিং-অগণনসাধ্য অপেক্ষকের অস্তিত্ব অনুসৃত হয়: $\Nat$ থেকে
~$\Nat$-এ সকল অপেক্ষকের সেট (আসলে শুধু $\Nat$ থেকে $\{0, 1\}$-এ সকল
অপেক্ষকের সেটও) !!{nonenumerable}, কিন্তু সব টুরিং যন্ত্রকে তালিকায়িত করা
যায় বলে টুরিং-গণনসাধ্য অপেক্ষকের সেটটি কেবল !!{denumerable}।

এমন \emph{নির্দিষ্ট} অপেক্ষক ও সমস্যাও সংজ্ঞায়িত করা যায়, যেগুলি যথাক্রমে
অগণনসাধ্য ও অনির্ণেয় বলে প্রমাণ করা যায়। তেমন একটি সমস্যা হলো তথাকথিত
\emph{থামার সমস্যা}। টুরিং যন্ত্রের নির্দেশগুলি তালিকায় লিখে সসীমভাবে তার
বর্ণনা দেওয়া যায়। টুরিং যন্ত্রের এমন একটি বর্ণনা, অর্থাৎ একটি টুরিং-
যন্ত্র প্রোগ্রাম, অবশ্যই অন্য টুরিং যন্ত্রের নিবেশ হিসেবে ব্যবহার করা যায়।
তাই অন্য টুরিং যন্ত্র সম্পর্কে প্রশ্ন নির্ণয় করে এমন টুরিং যন্ত্র বিবেচনা
করা যায়। বিশেষ আকর্ষণীয় একটি প্রশ্ন হলো: “দেওয়া টুরিং যন্ত্রটি নিবেশ
~$n$-এ শুরু করলে শেষ পর্যন্ত থামে কি?” এই প্রশ্ন নির্ণয় করতে পারে এমন
একটি টুরিং যন্ত্র থাকলে ভালো হতো: তাকে এমন একটি মান-নিয়ন্ত্রণকারী টুরিং
যন্ত্র ভাবা যায়, যা নিশ্চিত করে যে টুরিং যন্ত্রগুলি অসীম চক্র ইত্যাদিতে
আটকে পড়ে না। টুরিং যে আকর্ষণীয় তথ্যটি প্রমাণ করেছিলেন তা হলো, এমন কোনো
টুরিং যন্ত্র থাকতে পারে না। টুরিং যন্ত্র~$M$-এর একটি বর্ণনা ও কোনো সংখ্যা
~$n$ নিয়ে গঠিত নিবেশে শুরু করে, $M$ নিবেশ~$n$-এ শুরু করলে থামত কি না তার
উপর নির্ভর করে সর্বদা নির্গম $1$ অথবা~$0$ দিয়ে থামে—এমন একটিমাত্র টুরিং
যন্ত্র থাকতে পারে না।

নির্দিষ্ট অনির্ণেয় সমস্যার উদাহরণ পাওয়ার পর সেগুলির সাহায্যে অন্য
সমস্যাও অনির্ণেয় দেখানো যায়। যেমন, একটি বিখ্যাত অনির্ণেয় সমস্যা হলো,
“প্রথম-ক্রমের !!{formula}~$!A$ কি সিদ্ধ?” কোনো প্রথম-ক্রমের
!!{formula}~$!A$ নিবেশ হিসেবে নিয়ে, $!A$ সিদ্ধ কি না অনুযায়ী নির্গম
$1$ অথবা~$0$ দিয়ে থামবে বলে নিশ্চয়তা আছে—এমন কোনো টুরিং যন্ত্র নেই।
ঐতিহাসিকভাবে, এই সমস্যা কার্যকরভাবে সমাধান করার পদ্ধতি খোঁজার প্রশ্নটিকে
কেবল “সিদ্ধান্ত সমস্যা” বলা হতো; তাই আমরা বলি, সিদ্ধান্ত সমস্যা
অসমাধানযোগ্য। টুরিং ও চার্চ প্রায় একই সময়ে স্বাধীনভাবে এই ফল প্রমাণ
করেছিলেন বলে একে চার্চ--টুরিং উপপাদ্যও বলা হয়।

\end{document}
